\begin{proposition}\label{prop:decic95-excluded}
No elliptic curve over a number field of degree $10$ has a point
of order $95$.
\end{proposition}
\begin{proof}
Put $X=X_1(95)$ and $h=\langle11\rangle$.
We apply Lemma~\ref{lem:dyadic-cuspidal-diamond} with $d=10$.
For $f\leq7$, the Hasse interval contains no multiple of $95$.
The remaining possibilities are
\[
\begin{array}{c|c|c}
f&\#E(\mathbb F_{2^f})&a_{2^f}\\\hline
8&285&-28\\
9&475&38\\
10&1045&-20
\end{array}
\]
An ordinary trace in characteristic $2$ is odd, and none of the
displayed even traces is supersingular by
\cite[Theorem~4.1]{Waterhouse}.
Thus $Y_1(95)(\mathbb F_{2^f})$ is empty for $f\leq10$.
An exhaustive enumeration of twists over these ten finite fields
agrees with this conclusion. Independent integer finite-field
arithmetic verifies all $4122$ elliptic-curve point counts, for
both this proposition and Proposition~\ref{prop:decic111-excluded}.

The Tate labels give global cusp signature
\[
 1^{36},\quad4^9,\quad18^2,\quad36.
\]
The four types are defined over $\mathbb Q$,
$\mathbb Q(\zeta_5)$, $\mathbb Q(\zeta_{19})$ and
$\mathbb Q(\zeta_{95})^+$, respectively.
Since $2$ has orders $4$, $18$ and projective order $36$
in the corresponding cyclotomic groups, every global cusp orbit
remains a single closed orbit at $2$.
Every degree-$10$ cuspidal residue divisor therefore has a rational
cuspidal lift of degree $10$.

The projective diamond group is cyclic of order $36$, generated by
$\langle2\rangle$, and $11\equiv2^{12}\pmod{95}$.
Only characters of orders $9$, $18$ and $36$ are nontrivial on $h$.
Their relative cuspidal dimensions are all $8$, and their exact
winding dimensions are also all $8$.
The complete even-character computation has weighted dimension
$289=g(X)$; only the order-$12$ component has a non-full winding
submodule, and it is trivial on $h$.
An exported $8\times8$ modular Krylov matrix has full rank at $73$.
Consequently $(h-1)J_1(95)(\mathbb Q)$ is finite.

Finally, $h$ has order $3$ and
\[
 \operatorname{gon}_{\mathbb C}(X)>
 \frac{325\cdot4320}{32768}
 =\frac{43875}{1024}>40.
\]
Since $3\nmid10$, the lemma excludes the proposed divisor.
\end{proof}
